\documentclass[12pt,a4paper]{report}

% ── Encodage & langue ─────────────────────────────────────────
\usepackage[utf8]{inputenc}
\usepackage[T1]{fontenc}
\usepackage[french]{babel}

% ── Mise en page ──────────────────────────────────────────────
\usepackage[top=2.5cm, bottom=2.5cm, left=3cm, right=2.5cm]{geometry}
\usepackage{setspace}
\onehalfspacing

% ── Typographie & couleurs ────────────────────────────────────
\usepackage{xcolor}
\usepackage{lmodern}
\definecolor{gold}{RGB}{201,168,76}
\definecolor{darkbg}{RGB}{20,20,20}
\definecolor{graytext}{RGB}{100,95,88}

% ── Graphiques & tableaux ─────────────────────────────────────
\usepackage{graphicx}
\usepackage{booktabs}
\usepackage{tabularx}
\usepackage{longtable}
\usepackage{array}
\usepackage{multirow}

% ── Code source ───────────────────────────────────────────────
\usepackage{listings}
\lstset{
    basicstyle=\ttfamily\small,
    backgroundcolor=\color{gray!10},
    frame=single,
    rulecolor=\color{gray!40},
    breaklines=true,
    captionpos=b,
    keywordstyle=\color{blue!70},
    commentstyle=\color{green!60!black},
    stringstyle=\color{orange!80!black},
    numbers=left,
    numberstyle=\tiny\color{gray},
    stepnumber=1,
}

% ── Liens & PDF ───────────────────────────────────────────────
\usepackage{hyperref}
\hypersetup{
    colorlinks=true,
    linkcolor=gold,
    urlcolor=blue!70,
    citecolor=gold,
    pdftitle={VoyageIQ-Pro — Rapport de Projet},
    pdfauthor={Équipe VoyageIQ-Pro — ESTLC},
    pdfsubject={Système ERP de gestion de transport interurbain},
}

% ── Divers ────────────────────────────────────────────────────
\usepackage{enumitem}
\usepackage{fancyhdr}
\usepackage{titlesec}
\usepackage{tcolorbox}
\usepackage{amsmath}

% ── En-têtes & pieds de page ──────────────────────────────────
\pagestyle{fancy}
\fancyhf{}
\fancyhead[L]{\small\textcolor{graytext}{VoyageIQ-Pro}}
\fancyhead[R]{\small\textcolor{graytext}{ESTLC — 2025-2026}}
\fancyfoot[C]{\thepage}
\renewcommand{\headrulewidth}{0.4pt}

% ── Style des titres de chapitres ─────────────────────────────
\titleformat{\chapter}[display]
  {\normalfont\Large\bfseries\color{gold}}
  {\chaptertitlename\ \thechapter}{10pt}{\huge}

% ═════════════════════════════════════════════════════════════
\begin{document}

% ── Page de titre ─────────────────────────────────────────────
\begin{titlepage}
    \centering
    \vspace*{1.5cm}

    {\Large\textbf{ESTLC}}\\[0.3cm]
    {\large École Supérieure de Transport et de Logistique du Cameroun}\\[0.2cm]
    {\normalsize Département Transport \& Logistique}\\[1cm]

    \rule{\linewidth}{1.5pt}\\[0.4cm]
    {\Huge\bfseries\color{gold} VoyageIQ-Pro}\\[0.3cm]
    {\Large Système ERP de Gestion de Transport Interurbain}\\[0.2cm]
    \rule{\linewidth}{1.5pt}\\[1cm]

    {\large\textit{Rapport de Projet Professionnel}}\\[0.2cm]
    {\normalsize Année académique 2025--2026}\\[1.5cm]

    % Équipe
    \begin{minipage}{0.48\textwidth}
        \centering
        \textbf{Équipe Projet}\\[0.3cm]
        \begin{tabular}{ll}
            \toprule
            \textbf{Nom \& Prénom} & \textbf{Département} \\
            \midrule
            Membre 1 (Transport)   & Transport \& Logistique \\
            Membre 2 (Transport)   & Transport \& Logistique \\
            Membre 3 (Transport)   & Transport \& Logistique \\
            Membre 4 (Transport)   & Transport \& Logistique \\
            Membre 5 (Transport)   & Transport \& Logistique \\
            \textbf{Masse Masse Paul-Basthylle} & Informatique \\
            \bottomrule
        \end{tabular}
    \end{minipage}
    \hfill
    \begin{minipage}{0.45\textwidth}
        \centering
        \textbf{Encadrement}\\[0.3cm]
        \begin{tabular}{l}
            \toprule
            \textbf{Encadreur académique} \\
            \midrule
            Dr. Prosper \\
            Département Transport \& Logistique\\
            ESTLC \\
            \bottomrule
        \end{tabular}
    \end{minipage}

    \vfill
    {\normalsize Yaoundé, Cameroun --- \today}
\end{titlepage}

% ── Table des matières ────────────────────────────────────────
\tableofcontents
\newpage

% ─────────────────────────────────────────────────────────────
\chapter{Introduction Générale}
% ─────────────────────────────────────────────────────────────

\section{Contexte}

Le secteur du transport interurbain au Cameroun constitue un pilier essentiel de
l'économie nationale, assurant la mobilité des personnes et des biens entre les
principales agglomérations du pays. Malgré son importance stratégique, ce secteur
souffre encore d'un déficit notable en matière de digitalisation et de centralisation
des données opérationnelles.

La majorité des agences de transport gèrent encore leurs opérations quotidiennes
via des registres papier ou des tableurs non connectés, rendant impossible le suivi
en temps réel de la performance et la prise de décision rapide.

\section{Problématique}

Face à ce constat, la question centrale de notre projet est la suivante :

\begin{tcolorbox}[colback=gray!10, colframe=gold, title=Question centrale]
\textit{Comment concevoir et déployer une plateforme web intégrée permettant à
une agence de transport interurbain de centraliser ses données opérationnelles,
financières et humaines, tout en offrant des indicateurs de performance exploitables
en temps réel ?}
\end{tcolorbox}

\section{Objectifs du Projet}

\subsection{Objectif général}
Développer \textbf{VoyageIQ-Pro}, un système ERP web complet adapté aux besoins
spécifiques du transport interurbain camerounais.

\subsection{Objectifs spécifiques}
\begin{itemize}[leftmargin=1.5cm]
    \item Digitaliser la saisie quotidienne des rapports d'exploitation par ligne.
    \item Automatiser le suivi de la maintenance et des documents de la flotte.
    \item Analyser la rentabilité financière par ligne de transport.
    \item Intégrer un espace personnel pour les chauffeurs (courses, stats, localisation).
    \item Évaluer la satisfaction clientèle via un système d'avis et le NPS.
    \item Générer des rapports PDF et les diffuser par email et WhatsApp.
    \item Offrir une vitrine publique professionnelle pour l'agence.
\end{itemize}

\section{Périmètre et Limites}

Ce projet couvre la conception et le développement d'une application web accessible
via navigateur. La version actuelle (v2.0) ne comprend pas d'application mobile native,
bien que l'intégration d'une API REST permette une extension future dans ce sens.

% ─────────────────────────────────────────────────────────────
\chapter{Analyse du Système}
% ─────────────────────────────────────────────────────────────

\section{Identification des Acteurs}

Le système identifie trois grandes catégories d'acteurs :

\begin{table}[h!]
\centering
\caption{Acteurs du système VoyageIQ-Pro}
\begin{tabularx}{\textwidth}{llX}
\toprule
\textbf{Acteur} & \textbf{Espace} & \textbf{Rôle principal} \\
\midrule
Visiteur / Client   & Public      & Consultation de la vitrine, avis \\
Personnel (5 rôles) & Admin       & Gestion ERP complète selon niveau \\
Chauffeur           & Chauffeur   & Suivi personnel, courses, localisation \\
\bottomrule
\end{tabularx}
\end{table}

\section{Hiérarchie des Rôles Administratifs}

\begin{table}[h!]
\centering
\caption{Niveaux d'accès dans l'espace administratif}
\begin{tabular}{clll}
\toprule
\textbf{Niveau} & \textbf{Rôle} & \textbf{Accès} & \textbf{Auth} \\
\midrule
5 & Administrateur      & Complet (configuration, utilisateurs) & Téléphone + MDP \\
4 & Direction Générale  & Finance, analytique, KPI globaux      & Téléphone + MDP \\
3 & Chef d'Agence       & Flotte, saisie, rapports agence       & Téléphone + MDP \\
2 & Superviseur Terrain & Saisie des données, opérations        & Téléphone + MDP \\
1 & Auditeur            & Consultation rapports et alertes      & Téléphone + MDP \\
\bottomrule
\end{tabular}
\end{table}

\section{Modèle de Données}

Le système repose sur sept entités principales :

\begin{description}[leftmargin=2cm, style=nextline]
    \item[\texttt{Utilisateur}]
        Gère l'authentification par téléphone, le contrôle d'accès basé sur les
        rôles (RBAC), le profil enrichi (photo, matricule, coordonnées, WhatsApp).
    \item[\texttt{Chauffeur}]
        Acteur externe pouvant s'inscrire et accéder à un espace personnel.
        Suit un cycle de validation administrateur avant activation.
    \item[\texttt{CourseChauffeur}]
        Association Chauffeur $\leftrightarrow$ Saisie/Voyage, permettant le suivi
        individuel des performances de chaque chauffeur.
    \item[\texttt{Vehicule}]
        Suivi des caractéristiques techniques, kilométrage, échéances de
        maintenance, assurance et visite technique.
    \item[\texttt{Ligne}]
        Définition des itinéraires, paramètres tarifaires et fréquences de passage.
    \item[\texttt{Saisie}]
        Cœur du système. Enregistre les données quotidiennes : voyages, recettes,
        dépenses, incidents, qualité, NPS.
    \item[\texttt{Alerte}]
        Notifications proactives générées automatiquement par le moteur de règles.
\end{description}

% ─────────────────────────────────────────────────────────────
\chapter{Architecture et Implémentation}
% ─────────────────────────────────────────────────────────────

\section{Choix Technologiques}

\begin{table}[h!]
\centering
\caption{Stack technologique VoyageIQ-Pro}
\begin{tabularx}{\textwidth}{llX}
\toprule
\textbf{Couche} & \textbf{Technologie} & \textbf{Justification} \\
\midrule
Backend    & Python 3.10+ / Flask    & Micro-framework flexible, écosystème riche \\
ORM        & SQLAlchemy              & Abstraction BDD, migrations, sécurité \\
Base de données & SQLite (dev) / PostgreSQL (prod) & Légèreté en dev, robustesse en prod \\
Auth       & Flask-Login + Werkzeug  & Gestion sessions, hachage PBKDF2 \\
Sécurité   & Flask-WTF (CSRF)        & Protection formulaires \\
Frontend   & Jinja2 + CSS + JS       & Templates serveur, design custom \\
Rapports   & jsPDF / WeasyPrint      & Génération PDF côté client/serveur \\
Cartographie & Leaflet.js + OpenStreetMap & Localisation chauffeurs (open source) \\
\bottomrule
\end{tabularx}
\end{table}

\section{Structure du Projet}

\begin{lstlisting}[language=bash, caption=Arborescence principale VoyageIQ-Pro]
VoyageIQ-Pro/
├── app/
│   ├── blueprints/        # Modules routes (bp_<module>.py)
│   ├── models/            # Entités SQLAlchemy
│   ├── services/          # Logique métier (KPI, alertes)
│   ├── static/            # CSS, JS, images, fonts
│   │   ├── css/base/      # Variables, reset, animations
│   │   ├── css/components/# Sidebar, navbar, cards...
│   │   └── css/pages/     # Styles par page
│   ├── templates/
│   │   ├── admin/         # Pages espace administratif
│   │   ├── chauffeur/     # Pages espace chauffeur
│   │   ├── public/        # Vitrine publique
│   │   ├── base/          # Templates de base (héritage Jinja2)
│   │   └── components/    # Composants réutilisables
│   └── utils/             # Décorateurs, helpers
├── config/settings.py     # Configuration multi-environnements
├── database/seeds/        # Initialisation BDD
└── run.py                 # Point d'entrée
\end{lstlisting}

\section{Organisation des Blueprints}

Chaque module fonctionnel est encapsulé dans un fichier \texttt{bp\_<module>.py}
indépendant, enregistré dans \texttt{app/\_\_init\_\_.py} avec son préfixe d'URL.
Cette structure plate (sans sous-dossiers) améliore la lisibilité et simplifie
les imports.

\section{Services Métier}

\begin{description}[leftmargin=2cm]
    \item[\texttt{kpi\_service.py}]
        Centralise les calculs complexes : taux de remplissage, marge d'exploitation,
        consommation aux 100 km, NPS agrégé.
    \item[\texttt{alerte\_service.py}]
        Moteur de règles analysant l'état de la flotte pour générer automatiquement
        des alertes de maintenance, d'expiration de documents et d'anomalies financières.
\end{description}

% ─────────────────────────────────────────────────────────────
\chapter{Fonctionnalités Développées}
% ─────────────────────────────────────────────────────────────

\section{Espace Public (Vitrine)}
La page d'accueil présente l'agence avec des statistiques dynamiques,
les lignes desservies, les modules de la plateforme, un formulaire de contact
et les mentions légales. Un bouton \og Administration \fg{} donne accès
au login sans exposer l'interface interne.

\section{Authentification par Téléphone}
L'identifiant de connexion est le numéro de téléphone (format \texttt{+237XXXXXXXXX}
ou \texttt{XXXXXXXXX}). Le système accepte les deux formats de manière transparente.
Les chauffeurs disposent d'une page de connexion dédiée (\texttt{/chauffeur/login}).

\section{Espace Administratif}
Dix modules accessibles selon le niveau hiérarchique :
Dashboard, Saisie, Flotte, Lignes, Finance, Opérations, Clientèle, Analytique,
Alertes, Rapports. Les admins accèdent en plus à la gestion des utilisateurs et chauffeurs.

\section{Espace Chauffeur}
Les chauffeurs accèdent à un espace personnel après inscription et validation :
dashboard personnel, historique des courses, statistiques de performance,
suivi des maintenances, carte de localisation en temps réel, et modification du profil.

\section{Génération de Rapports}
Les rapports (exploitation journalière, finance mensuelle, bilan flotte) sont
générables en PDF et envoyables par email et/ou WhatsApp selon les préférences
de notification de chaque utilisateur.

% ─────────────────────────────────────────────────────────────
\chapter{Gestion de la Sécurité}
% ─────────────────────────────────────────────────────────────

La sécurité est articulée autour de quatre axes :

\begin{enumerate}[leftmargin=1.5cm]
    \item \textbf{Hachage des mots de passe} : Algorithme PBKDF2 via Werkzeug.
    \item \textbf{Protection CSRF} : Intégrée à tous les formulaires via Flask-WTF.
    \item \textbf{Contrôle d'accès par rôle (RBAC)} : Décorateurs Flask restreignant
          les routes selon le niveau hiérarchique (\texttt{niveau 1} à \texttt{5}).
    \item \textbf{Séparation des sessions} : Les Chauffeurs et les Utilisateurs admin
          utilisent des préfixes d'identifiant distincts dans Flask-Login
          (\texttt{c-\{id\}} pour les chauffeurs).
\end{enumerate}

% ─────────────────────────────────────────────────────────────
\chapter{Perspectives d'Évolution}
% ─────────────────────────────────────────────────────────────

\begin{itemize}[leftmargin=1.5cm]
    \item \textbf{Application mobile} pour les chauffeurs (iOS / Android) via l'API REST.
    \item \textbf{Géolocalisation en temps réel} des véhicules (WebSocket + Leaflet.js).
    \item \textbf{Intégration Mobile Money} (MTN MoMo, Orange Money) pour les recettes digitales.
    \item \textbf{Module RH} : gestion des congés, contrats et fiches de paie des chauffeurs.
    \item \textbf{Tableau de bord exécutif} avec export Excel et rapports automatisés périodiques.
    \item \textbf{Migration PostgreSQL} pour l'environnement de production multi-utilisateurs.
\end{itemize}

% ─────────────────────────────────────────────────────────────
\chapter{Conclusion}
% ─────────────────────────────────────────────────────────────

\textbf{VoyageIQ-Pro} constitue une réponse concrète et opérationnelle aux défis
de digitalisation du transport interurbain camerounais. En centralisant les données
opérationnelles, financières et humaines dans une seule plateforme, le système
permet aux décideurs d'agir sur la base d'informations fiables et en temps réel.

La conception modulaire (blueprints, services découplés, modèles enrichis) garantit
une maintenabilité optimale et une capacité d'évolution progressive. L'intégration
d'un espace dédié aux chauffeurs, d'un système de notifications multicanal et d'une
vitrine publique professionnelle place VoyageIQ-Pro au-delà d'un simple outil de
saisie, en faisant une véritable plateforme de pilotage d'agence.

Ce projet, mené en équipe pluridisciplinaire (Transport \& Informatique) sous
l'encadrement du Dr.~Prosper, illustre la complémentarité entre expertise métier
et compétences techniques dans la construction de solutions numériques adaptées
au contexte africain.

% ─────────────────────────────────────────────────────────────
% Annexes
% ─────────────────────────────────────────────────────────────
\appendix

\chapter{Comptes de Démonstration}

\begin{table}[h!]
\centering
\caption{Accès de démonstration — Espace Admin}
\begin{tabular}{lll}
\toprule
\textbf{Téléphone} & \textbf{Mot de passe} & \textbf{Rôle} \\
\midrule
+237690000001 & Admin@VIQ2026  & Administrateur \\
+237677000010 & DG@VIQ2026     & Direction Générale (Mbarga J-P) \\
+237655000100 & Chef1@VIQ2026  & Chef d'Agence Yaoundé (Biyong P.) \\
+237666000200 & Chef2@VIQ2026  & Chef d'Agence Douala (Ngom B.) \\
+237655001001 & Sup1@VIQ2026   & Superviseur Yaoundé (Ateba R.) \\
+237677002001 & Audit1@VIQ2026 & Auditeur (Nkolo F.) \\
\bottomrule
\end{tabular}
\end{table}

\begin{table}[h!]
\centering
\caption{Accès de démonstration — Espace Chauffeur}
\begin{tabular}{lll}
\toprule
\textbf{Téléphone} & \textbf{Mot de passe} & \textbf{Chauffeur} \\
\midrule
+237677100001 & Chauf1@VIQ2026 & Tsafack Hervé (validé) \\
+237655100002 & Chauf2@VIQ2026 & Nganou Alphonse (validé) \\
+237699100003 & Chauf3@VIQ2026 & Kamga Rodrigue (validé) \\
+237677200004 & Fouda@2026     & Fouda Serge (en attente) \\
\bottomrule
\end{tabular}
\end{table}

\chapter{Lignes de Transport Configurées}

\begin{table}[h!]
\centering
\caption{Lignes de démonstration}
\begin{tabular}{llrrl}
\toprule
\textbf{Code} & \textbf{Trajet} & \textbf{Km} & \textbf{Tarif (FCFA)} & \textbf{Fréq./jour} \\
\midrule
L01 & Yaoundé — Douala    & 250 & 3\,500 & 8 \\
L02 & Yaoundé — Bafoussam & 310 & 4\,000 & 6 \\
L03 & Douala — Kribi      & 195 & 2\,500 & 5 \\
L04 & Yaoundé — Bertoua   & 350 & 4\,500 & 4 \\
L05 & Douala — Bafoussam  & 375 & 4\,500 & 4 \\
L06 & Yaoundé — Ebolowa   & 160 & 2\,000 & 4 \\
L07 & Douala — Limbé      &  70 & 1\,000 & 10 \\
\bottomrule
\end{tabular}
\end{table}

\end{document}
